%%%%%%%%%%  scan idx 278  |  folio 267  %%%%%%%%%%

Recordemos, en forma de lema, el resultado siguiente de EGA
$\mathrm{Err}_{\mathrm{IV}}$ 53, 21.4.13.

\medskip
\noindent\underline{Lema 4.5.} \underline{Sea
$f:X\longrightarrow Y$ un morfismo liso y suprayectivo con fibras
íntegras, donde $Y$ es normal e íntegro y tiene punto genérico $y$.
Sea $\mathbb L$ un módulo invertible sobre $X$ cuya restricción a
$X_y$ es trivial. Entonces existe un módulo invertible $\mathbb M$
sobre $Y$ tal que}
\[
 \mathbb L\simeq f^*\mathbb M.
\]

\medskip
\noindent\underline{Proposición 4.6.} \underline{Sean $P,Q$ dos
grupos algebraicos lisos y conexos sobre el cuerpo $k$, con $P$ afín.
Supongamos además, bien que $P$ admite una serie de composición cuyos
factores son toros o son isomorfos al grupo aditivo $\underline G_a$
(condición que se cumple en particular cuando $k$ es perfecto [SGA 3,
XVII 7.2.1, 4.1.1 (i)$\Longleftrightarrow$(ii), y 4.1.5]), bien que
$Q$ es una extensión de una variedad abeliana $B$ por un grupo
algebraico afín liso y conexo que satisface la condición precedente.
Bajo cualquiera de estas hipótesis, todo módulo invertible
birigidificado $\mathbb L$ sobre $P\times Q$ es trivial; con mayor
razón, por 4.3, toda biextensión de $(P,Q)$ por $\underline G_m$ es
trivial.}

\noindent\underline{Demostración.} Observemos primero que todo
automorfismo de un módulo invertible birigidificado es la identidad
(4.3). La teoría del descenso muestra, por tanto, que basta demostrar
que $\mathbb L$ se hace trivial tras una extensión finita separable de
$k$. El paso al límite nos reduce inmediatamente al caso en que $k$
es separablemente cerrado.

Todo toro sobre un cuerpo tal es escindido. Por ello, si $P$ admite
una serie de composición como en 4.6, admite también una cuyos
cocientes son isomorfos a $\underline G_a$ o $\underline G_m$. La
aplicación repetida de 4.5, junto con
\[
 \Pic(\underline G_{a_K})=\Pic(\underline G_{m_K})=0
\]
para todo cuerpo $K$, da
\[
 \Pic(P)=0.
\]

%%%%%%%%%%  scan idx 279  |  folio 268  %%%%%%%%%%

Esta igualdad sigue siendo válida después de toda extensión del cuerpo
de base $k$. Aplicando 4.5 a
\[
 \pr_2:P\times Q\longrightarrow Q
\]
y a $\mathbb L$, encontramos un módulo invertible $\mathbb M$ sobre
$Q$ tal que $\mathbb L\simeq\pr_2^*\mathbb M$. Como la restricción de
$\mathbb L$ a $e_P\times Q$ es trivial, $\mathbb M$ es trivial, y por
tanto también $\mathbb L$. Esto prueba el primer caso.

En el segundo caso, sea $Q$ una extensión de un esquema abeliano $B$
por un grupo $M$ que admite una serie de composición con cocientes
isomorfos a $\underline G_m$ o $\underline G_a$. Apliquemos 4.5 a
\[
 f:X=P\times Q\longrightarrow Y=P\times B.
\]
La fibra sobre el punto genérico $y$ de $P\times B$ es un
$M_y$-torsor. Es necesariamente trivial por la estructura particular
de $M$ y por las igualdades clásicas
\[
 H^1(K,\underline G_a)=0,
 \qquad
 H^1(K,\underline G_m)=0,
\]
válidas para todo cuerpo $K$, aquí $K=k(y)$. Como $\Pic(M)=0$, se
tiene $\Pic(X_y)=0$; por tanto, 4.5 se aplica y muestra que $\mathbb L$
es la imagen inversa de un módulo invertible $\mathbb N$ sobre
$Y=P\times B$.

Sean $\mathbb N_1$ y $\mathbb N_2$ las restricciones de $\mathbb N$
a $P\times e_B$ y $e_P\times B$, respectivamente. Reemplazando
$\mathbb N$ por
\[
 \mathbb N\otimes\pr_1^*(\mathbb N_1)^{-1}
 \otimes\pr_2^*(\mathbb N_2)^{-1},
\]
donde $\pr_1$ y $\pr_2$ son las dos proyecciones de
$Y=P\times\sisi{A}{B}$, podemos suponer que $\mathbb N$ está
birigidificado y que $\mathbb L$ es su imagen inversa
birigidificada. Quedamos así reducidos al caso en que el propio $Q$ es
un esquema abeliano.

Entonces $\mathbb L$ define un morfismo canónico
\begin{equation}\tag{*}
 P\longrightarrow\underline{\Pic}_{Q/k}
\end{equation}
que lleva la identidad en la identidad. Como $P$ es conexo, este
morfismo se factoriza por el esquema abeliano dual

%%%%%%%%%%  scan idx 280  |  folio 269  %%%%%%%%%%

\[
 Q'=\underline{\Pic}^{0}_{Q/k}.
\]
Es bien sabido que todo morfismo de esquemas de un grupo algebraico
afín, liso y conexo a un esquema abeliano es constante. Es una
afirmación geométrica, de modo que podemos suponer que $k$ es
algebraicamente cerrado; un dévissage la reduce entonces
inmediatamente a $P=\underline G_a$ o $P=\underline G_m$, y toda
aplicación racional de $\mathbb P^1_k$ a un esquema abeliano sobre $k$
es constante.

La constancia, y por tanto la nulidad, de (*) significa que
$\mathbb L$ es la imagen inversa de un módulo invertible $\mathbb M$
sobre $P$. Como $\mathbb L$ es trivial sobre
$P\times e_{\sisi{A}{Q}}$, $\mathbb M$ es trivial, y por consiguiente
$\mathbb L$ también lo es. Esto termina la demostración.
